\begin{frontmatter}

\title{Screening for Bid Rigging with Frequent Losers in Public Procurement\texorpdfstring{\footnote{This version: April 2026.}}{}}

\author[insper]{Darcio Genicolo-Martins\corref{cor1}}
\ead{darciogm1@insper.edu.br}
\author[insper]{Paulo Furquim de Azevedo}
\affiliation[insper]{organization={INSPER},
  city={S\~ao Paulo}, country={Brazil}}
\cortext[cor1]{Corresponding author.}

\begin{abstract}
Bid-rigging cartels need to simulate competition, and a common device
is deploying firms that submit deliberately losing bids. We turn this
footprint into a firm-level diagnostic tool: \emph{frequent losers}
(FL)---firms that never win yet participate in abnormally many
tenders. In within-item comparisons across 4.5 million tender-items
from Brazilian public procurement (2009--2019), FL presence is
associated with a 3.6--7.7\% conditional price gap. Among 2,735 FL
firms, 7.1\% co-participate with convicted cartelists---3.5 times the
baseline rate. Multiple diagnostics produce results coherent with cover bidding as
one plausible interpretation. The screen requires only participation records
and can operate where bid-level tools cannot; a flag is a starting
point for investigation, not a verdict.
\end{abstract}

\end{frontmatter}

\noindent\textbf{Keywords:} bid rigging, cover bidding, public procurement, cartel screening, frequent losers

\noindent\textbf{JEL No.:} D44, D73, H57, K21, L41

\bigskip
